Bahan dan peralatan yang digunakan dalam penelitian ini meliputi seluruh komponen yang mendukung proses penelitian secara menyeluruh, baik pada tahap simulasi algoritma, analisis normalisasi \emph{step-size}, maupun implementasi pada FPGA. Bahan yang digunakan mencakup sinyal uji dan referensi ilmiah, sedangkan peralatan yang digunakan berupa perangkat keras dan perangkat lunak yang menunjang proses akuisisi sinyal, simulasi, sintesis, pengukuran, dan analisis data. Rincian bahan dan peralatan tersebut disajikan pada Tabel~\ref{tab:bahan_peralatan}.

\begingroup
\small
\renewcommand{\arraystretch}{1.3}
\begin{longtable}{|>{\centering\arraybackslash}p{0.8cm}|>{\raggedright\arraybackslash}p{3.3cm}|>{\raggedright\arraybackslash}p{5.2cm}|>{\raggedright\arraybackslash}p{4.6cm}|}
  \caption{Bahan dan Peralatan yang Digunakan}
  \label{tab:bahan_peralatan} \\
  \hline
  \textbf{No.} & \textbf{Nama Peralatan / Komponen} & \textbf{Spesifikasi} & \textbf{Fungsi} \\
  \hline
  \endfirsthead
  \caption[]{Bahan dan Peralatan yang Digunakan (lanjutan)} \\
  \hline
  \textbf{No.} & \textbf{Nama Peralatan / Komponen} & \textbf{Spesifikasi} & \textbf{Fungsi} \\
  \hline
  \endhead
  1. & \emph{Board} FPGA & Digilent Arty A7-35T, Xilinx Artix-7 XC7A35T, 20.800 LUT, 41.600 FF, 90 DSP48E1, 100 MHz & Platform implementasi algoritma FXLMS dan FXNLMS \\
  \hline
  2. & Mikrofon Referensi & INMP441, MEMS, antarmuka I2S, 24-bit, \emph{omnidirectional}, 1 unit & Akuisisi sinyal kebisingan di luar \emph{earcup} \\
  \hline
  3. & Mikrofon \emph{Error} & INMP441, MEMS, antarmuka I2S, 24-bit, \emph{omnidirectional}, 1 unit & Akuisisi sinyal residu di dalam \emph{earcup}, dekat telinga \\
  \hline
  4. & DAC Audio & PCM5102A, antarmuka I2S, 32-bit, 112 dB SNR & Konversi sinyal \emph{anti-noise} digital menjadi analog \\
  \hline
  5. & Penguat \emph{Headphone} & TPA6132A2, stereo, \emph{class}-G, 2,3-5,5 V & Penguatan sinyal analog untuk menggerakkan \emph{speaker} \\
  \hline
  6. & \emph{Speaker Driver} & Dinamis 40 mm, 32 $\Omega$, 2 unit (1 aktif, 1 cadangan) & Pembangkit sinyal \emph{anti-noise} \\
  \hline
  7. & Regulator Tegangan Rendah \emph{Noise} & LP5907, LDO 3,3 V, 250 mA, \emph{noise} 6,5 $\mu$Vrms & Catu daya mikrofon bebas riak \\
  \hline
  8. & Mikrofon Ukur & MIC-L20 \emph{Clip On Microphone} USB & Pengukuran independen tingkat peredaman kebisingan (dB) \\
  \hline
  9. & Sumber Kebisingan & \emph{Speaker} aktif/\emph{studio monitor} & Pembangkit \emph{white noise} dan \emph{pink noise} uji \\
  \hline
  10. & Komputer / Laptop & Windows 11/Linux & Menjalankan perangkat lunak sintesis HDL, simulasi, dan analisis data \\
  \hline
  11. & Vivado Design Suite & Perangkat lunak sintesis dan implementasi FPGA Xilinx & Sintesis HDL serta pengambilan laporan utilisasi LUT, \emph{Flip-Flop}, memori, dan DSP \emph{Slices} \\
  \hline
  12. & MATLAB / Python & Perangkat lunak simulasi dan analisis numerik & Simulasi sistem ANC, implementasi FXLMS dan FXNLMS, serta perhitungan metrik kinerja \\
  \hline
  13. & Sinyal Uji & \emph{White noise}, \emph{pink noise}, dan rekaman bising lingkungan nyata & Sinyal kebisingan referensi untuk simulasi dan pengujian prototipe \\
  \hline
  14. & Referensi jurnal dan artikel ilmiah & Jurnal, prosiding, dan artikel terkait ANC, FXLMS, FXNLMS, dan implementasi FPGA & Dasar teoritis dalam pemilihan algoritma, parameter evaluasi, dan implementasi pada FPGA \\
  \hline
\end{longtable}
\endgroup
